% !TeX root = ../../../main.tex
\subsection{UMLS Mapping for Diagnosis Standardization}
\label{ch05:apd:umls-mapping}

To standardize the diagnoses extracted by LLMs from the case reports, we mapped the free-text diagnoses to concepts in the Unified Medical Language System (UMLS). 
This process involved the following steps:

\begin{enumerate}
    \item \textbf{Initial Extraction}: For each case report, we prompted the LLM to extract diagnoses. 
    This yielded a list of free-text diagnoses for each case report.

    \item \textbf{Entity Linking}: We used ScispaCy's entity linker model 
    \texttt{(en\_core\_sci\_lg)} with the UMLS knowledge base to map each extracted diagnosis to corresponding UMLS Concept Unique Identifiers (CUIs). 
    \ChFiveRebuttal{Each complete extracted diagnosis string is now supplied directly to the UMLS candidate generator.}

    \item \textbf{Filtering and Disambiguation}: For each diagnosis, we:
    \begin{itemize}
        \item Retained only mappings with a confidence score above 0.85.
        \item \ChFiveRebuttal{When candidates are tied or near-tied, the selection procedure favors exact normalized canonical-name matches, recognized canonical suffixes, and concepts whose canonical name contains the complete diagnosis phrase.}
    \end{itemize}

    \item \textbf{Concept Normalization}: After mapping to CUIs, we:
    \begin{itemize}
        \item Retrieved the canonical names from UMLS Metathesaurus.
        % \item Applied semantic type filtering to ensure mapped concepts were of types relevant to diagnoses (e.g., "Disease or Syndrome", "Pathologic Function", "Mental or Behavioral Dysfunction")
        % \item Grouped synonym concepts where appropriate to reduce redundancy
    \end{itemize}
\end{enumerate}

% \input{Scientific_data/Supplementary/US_disease_prevalence}

This mapping process allowed us to standardize the diverse expressions of diagnoses extracted by LLMs into a consistent terminology system, enabling more reliable frequency analysis and co-occurrence patterns. 
% \textcolor{red}{Jeremy: Quantify the mapping success rate. Any effects of confidence score or justification for threshold? @Shariar}
Across the full PMOA-TTS corpus, LLM-extracted diagnoses yielded
$N_{\text{raw}} = 63{,}715$ unique free-text diagnosis strings.
Of these, $N_{\text{mapped}} = 41{,}216$ ($64.7\%$) were successfully mapped to at least one
UMLS CUI with a confidence score $\ge 0.85$.
When weighted by frequency of occurrence, mappings at this threshold covered
$91.7\%$ of all diagnosis mentions, indicating that commonly occurring clinical
diagnoses were more likely to receive high-confidence mappings.
Mapping coverage decreased monotonically with increasing confidence threshold.
We therefore selected a threshold of 0.85 as a conservative operating point for
downstream normalization, providing broad coverage of frequently occurring diagnoses
while maintaining a high-confidence mapping criterion.


